% A TVCG journal paper typeset from the Building Occupancy notebook. Every
% number, table row, and figure reads the notebook's own results, so the paper
% follows its scope and threshold controls. DESIGN.md records the typographic
% decisions.
\documentclass[journal]{vgtc}

\onlineid{0}
\vgtccategory{Research}
\nocopyrightspace
\renewcommand{\manuscriptnotetxt}{}

\usepackage{booktabs}
\usepackage{newtxmath}
\usepackage{siunitx}
\usepackage[en-GB, calc]{datetime2}
\usepackage{marimo}

% vgtc looks for EPS figures when the engine is not pdfTeX, and Tectonic runs
% XeTeX, so name the formats that the notebook's figures arrive in.
\DeclareGraphicsExtensions{.pdf,.png,.jpg,.jpeg}

% The text reflows with every scope and threshold, so keep paragraphs whole at
% column and page breaks.
\widowpenalty=10000
\clubpenalty=10000
\brokenpenalty=10000
\hyphenation{Feldheim vacant}

% Numbers take the face around them, so the Helvetica abstract and captions
% print Helvetica digits, and a percent sign closes up to its number.
\sisetup{mode=match, reset-text-family=false, reset-text-series=false,
  reset-text-shape=false, group-separator={,}, group-minimum-digits=4,
  quantity-product={}}
\DeclareMarimoFormat{count}{round-mode=places, round-precision=0}
\DeclareMarimoFormat{reading}{round-mode=places, round-precision=1}
\DeclareMarimoFormat{pct}{scale=100, round-mode=places, round-precision=1}
\DeclareMarimoFormat{share}{scale=100, round-mode=places, round-precision=1,
  unit=\percent}
\DeclareMarimoFormat{weight}{round-mode=places, round-precision=2}

\DTMsettimestyle{default}
\DTMsetup{showseconds=false}
\DTMnewdatestyle{day}{\renewcommand*\DTMdisplaydate[4]{%
  \DTMweekdayname{##4} \number##3\ \DTMmonthname{##2}}%
  \renewcommand*\DTMDisplaydate{\DTMdisplaydate}}
\DTMnewdatestyle{longday}{\renewcommand*\DTMdisplaydate[4]{%
  \DTMweekdayname{##4} \number##3\ \DTMmonthname{##2} \number##1}%
  \renewcommand*\DTMDisplaydate{\DTMdisplaydate}}
\DTMnewdatestyle{shortday}{\renewcommand*\DTMdisplaydate[4]{%
  \DTMshortweekdayname{##4} \number##3\ \DTMshortmonthname{##2}}%
  \renewcommand*\DTMDisplaydate{\DTMdisplaydate}}

\newcommand{\cotwo}{\texorpdfstring{CO\textsubscript{2}}{CO₂}}
\newcommand{\period}{\marimodate[day]{scope_summary.period_start} to
  \marimodate[longday]{scope_summary.period_end}}
\newcommand{\threshold}{\marimonum[weight]{model_summary.threshold}}
\newcommand{\ifoccupied}[2]{\IfMarimoTF{scope_summary.occupied}{#1}{#2}}

% The paper names the notebook's scopes and sensors by their keys, so the
% notebook's control labels and column labels stay in the notebook. A key
% without a name prints as it is, and each read here carries a plain-text
% fallback, because \csname expands it.
\newcommand{\paperterm}[2]{%
  \ifcsname #1:#2\endcsname \csname #1:#2\endcsname \else #2\fi}
\expandafter\def\csname scope:all\endcsname{across the whole record}
\expandafter\def\csname scope:operating-hours\endcsname{on weekdays from 07:00
  to 19:00}
\expandafter\def\csname scope:off-hours\endcsname{outside weekday working hours}
\expandafter\def\csname sensor:CO2\endcsname{\cotwo{}}
\expandafter\def\csname sensor:Light\endcsname{light}
\expandafter\def\csname sensor:Temperature\endcsname{temperature}
\expandafter\def\csname sensor:Humidity\endcsname{humidity}
\expandafter\def\csname unit:CO2\endcsname{\cotwo{} (ppm)}
\expandafter\def\csname unit:Light\endcsname{Light (lux)}
\expandafter\def\csname unit:Temperature\endcsname{Temperature (\si{\degreeCelsius})}
\expandafter\def\csname unit:Humidity\endcsname{Humidity (\%)}
\newcommand{\scope}{\paperterm{scope}{\marimovalue[all]{analysis_scope.value}}}

% vgtc sets a caption outside a float with the caption package, which drops
% its text under XeTeX. The teaser numbers and sets its own caption instead.
\newcommand{\teasercaption}[2]{%
  \par\vspace{4pt}\refstepcounter{figure}\label{#1}%
  {\scriptsize\sffamily\leftskip=0pt\rightskip=0pt\parfillskip=0pt plus 1fil
   Figure~\thefigure: #2\par}}

\title{A Transparent Light and \cotwo{} Score for Room Occupancy}

\author{Building Occupancy Notebook}
\authorfooter{
  \item The Building Occupancy notebook computes every number, table, and
  figure in this paper, and marimo Studio typesets it with \LaTeX. The source
  is at \url{https://github.com/marimo-team/marimo-studio}.
}

\abstract{%
  Ambient sensors already sit in most offices, yet occupancy models built on
  them are rarely read by the people who run the building. We study
  \marimonum{scope_summary.observations} readings of temperature, humidity,
  light, and \cotwo{} recorded in Room~01 \scope{}, from \period.
  \ifoccupied{%
    The room is occupied in \marimonum[share]{scope_summary.occupancy_rate}
    of them. A score that adds normalized light and \cotwo{} with hand-set
    weights of \marimonum[weight]{occupancy_parameters.model.light_weight} and
    \marimonum[weight]{occupancy_parameters.model.co2_weight} labels
    \marimonum[share]{model_summary.accuracy} of the readings correctly in
    sample, with \marimonum[share]{model_summary.precision} precision and
    \marimonum[share]{model_summary.recall} recall at a threshold of
    \threshold.%
  }{%
    The room is never occupied in these readings, so the same light and
    \cotwo{} score is measured by its false alarms: it labels
    \marimonum[share]{model_summary.accuracy} of the readings correctly at a
    threshold of \threshold.%
  }
  The notebook that computes each number, table row, and figure also controls
  the scope and the threshold, so the paper changes with them.
}

\keywords{Occupancy detection, ambient sensing, interpretable models, building
analytics, reproducible documents.}

\teaser{
  \centering
  \marimographics[width=\linewidth, alt={Carbon dioxide and light in Room 01
  over the observation period, with occupied intervals shaded.}]{room_timeline}
  \teasercaption{fig:teaser}{%
    \cotwo{} (top) and light (bottom) in Room~01, averaged over ten-minute
    intervals.
    \ifoccupied{Shaded intervals are mostly occupied. Both signals rise within
    minutes of arrival and fall after departure.}{No interval is occupied, so
    the signals show the building's schedule and daylight without anyone in
    the room.}%
  }
}

\begin{document}

\firstsection{Introduction}

\maketitle

Lighting, heating, and ventilation run on schedules, while rooms empty and
fill on their own. Occupancy estimates close that gap, and the sensors they
need are often installed already. Candanedo and
Feldheim~\cite{Candanedo2016AccurateOccupancyDetection} showed that
statistical learning models detect the occupancy of an office from light,
temperature, humidity, and carbon dioxide (\cotwo{}) with high accuracy. A
facilities team, however, has to trust a rule before it switches equipment
off, and a rule reads best when every term in it can be pointed at on a chart.

This paper reads readings from that study through a deliberately simple
score. Its contributions are:
\begin{itemize}
  \item \ifoccupied{a profile of when Room~01 is used}{a profile of the hours
    the readings cover} (\cref{fig:teaser,fig:heatmap});
  \item \ifoccupied{a comparison of each sensor while the room is occupied and
    vacant}{each sensor's baseline while the room is vacant}
    (\cref{tab:sensors,fig:distributions});
  \item a two-signal score and the tradeoff its threshold controls
    (\cref{fig:threshold,tab:confusion}).
\end{itemize}

A marimo notebook~\cite{marimo} computes the values and draws the figures with
Matplotlib~\cite{Hunter2007Matplotlib}, and marimo Studio typesets this
\LaTeX{} source with Tectonic~\cite{Tectonic} whenever they change.

\begin{figure}[tb]
  \centering
  \marimographics[width=\columnwidth, alt={A grid of hours by weekdays shaded
  by the share of occupied readings.}]{occupancy_heatmap}
  \caption{Share of occupied readings by hour and weekday. White cells hold no
  readings.}
  \label{fig:heatmap}
\end{figure}

\begin{figure*}[tb]
  \centering
  \marimographics[width=\textwidth, alt={Four histograms comparing carbon
  dioxide, light, temperature, and humidity while the room is vacant and
  occupied.}]{sensor_distributions}
  \caption{Distribution of each sensor while Room~01 is
  \ifoccupied{vacant (gray) and occupied (orange), each scaled to its own
  peak, so the shapes compare across states of different size}{vacant, scaled
  to its peak}.}
  \label{fig:distributions}
\end{figure*}

\section{Room Observations}
\label{sec:data}

The readings come from the training split of the UCI Occupancy Detection
dataset~\cite{Candanedo2016OccupancyDetectionDataset}. Sensors in Room~01
logged temperature, relative humidity, light, and \cotwo{} once a minute, and
time-stamped pictures of the room record whether it was occupied. This paper
reads the \marimonum{scope_summary.observations} readings taken \scope{}.
\ifoccupied{%
  The room is occupied in \marimonum{scope_summary.occupied} of them,
  \marimonum[share]{scope_summary.occupancy_rate}, or about
  \marimonum[reading]{scope_summary.estimated_occupied_hours}~hours.
  Occupancy peaks on \marimodate[day]{daily_occupancy_peak.day}, when the room
  is in use for \marimonum[share]{daily_occupancy_peak.occupancy_rate} of the
  day's readings. \Cref{fig:heatmap} shows that use follows the working day:
  it starts in the morning, pauses at midday, and ends in the early evening.%
}{%
  None of these readings is occupied, and \cref{fig:heatmap} shows the hours
  they cover.%
}
\Cref{tab:days} lists each day with its share of occupied readings and its
\cotwo{} levels.

\begin{table}[tb]
  \caption{Readings per day, with the share that is occupied and the mean and
  peak \cotwo{} concentration.}
  \label{tab:days}
  \scriptsize%
  \centering%
  \begin{tabular}{@{}l r r r r@{}}
    \toprule
    & & & \multicolumn{2}{c@{}}{\cotwo{} (ppm)} \\
    \cmidrule(l){4-5}
    Day & Readings & Occupied (\%) & Mean & Peak \\
    \midrule
    \marimorows{daily_room_profile}{%
      \marimodate[shortday]{#1.day}
      & \marimonum{#1.observations}
      & \marimonum[pct]{#1.occupancy_rate}
      & \marimonum[count]{#1.mean_co2}
      & \marimonum[count]{#1.peak_co2} \\}
    \bottomrule
  \end{tabular}%
\end{table}

\section{Sensor Conditions}
\label{sec:sensors}

\Cref{tab:sensors} compares the mean of each sensor while the room is
\ifoccupied{vacant and occupied}{vacant}, and \cref{fig:distributions} compares
their distributions.
\IfMarimoTF{sensor_separation_peak}{%
  The means of
  \paperterm{sensor}{\marimovalue[?]{sensor_separation_peak.key}} lie furthest
  apart, \marimonum[share]{sensor_separation_peak.relative_separation} of its
  range. Light and \cotwo{} respond to people directly: the lights switch on
  with them, and their breathing raises \cotwo{} while they stay. Temperature
  also rises with occupancy, but it follows the heating and the weather as
  well, so the score in \cref{sec:score} reads light and \cotwo{}.%
}{%
  Without occupied readings, the readings cannot separate the states, and the
  table shows each sensor's vacant baseline.%
}

\begin{table}[tb]
  \caption{Mean sensor readings while Room~01 is vacant and occupied.
  \ifoccupied{Separation is the gap between the means as a share of the
  sensor's range.}{No reading in this scope is occupied, so the occupied means
  and their separation are empty.}}
  \label{tab:sensors}
  \scriptsize%
  \centering%
  \begin{tabular}{@{}l r r r@{}}
    \toprule
    Sensor & Vacant & Occupied & Separation (\%) \\
    \midrule
    \marimorows{sensor_profiles}{%
      \paperterm{unit}{\marimovalue[?]{#1.key}}
      & \marimonum[reading]{#1.vacant}
      & \marimonum[reading]{#1.occupied}
      & \marimonum[pct]{#1.relative_separation} \\}
    \bottomrule
  \end{tabular}%
\end{table}

\section{A Transparent Score}
\label{sec:score}

The score combines the two sensors that respond to people most directly.
Light $\ell$ is divided by its
\marimonum[count, scale=100]{occupancy_parameters.model.normalization_quantile}th
percentile $\ell_q$ and clipped to $[0, 1]$, and \cotwo{} $c$ is scaled from
its minimum $c_0$ to the same percentile $c_q$. \Cref{eq:score} adds the two
with weights set by hand:
\begin{equation}
  \label{eq:score}
  s = \marimonum[weight]{occupancy_parameters.model.light_weight}\,
      \frac{\min(\ell, \ell_q)}{\ell_q}
    + \marimonum[weight]{occupancy_parameters.model.co2_weight}\,
      \frac{\min(\max(c, c_0), c_q) - c_0}{c_q - c_0}.
\end{equation}
A reading counts as occupied when $s$ reaches the threshold $\tau$. Each term
is a sensor a facilities team already reads, and each weight states how much
that sensor counts.

\section{Results}
\label{sec:results}

\begin{figure}[tb]
  \centering
  \marimographics[width=\columnwidth, alt={Accuracy, precision, and recall
  across thresholds above the distribution of scores while vacant and
  occupied.}]{threshold_figure}
  \caption{Top: \ifoccupied{accuracy (solid), precision (dashed), and recall
  (dash-dotted)}{accuracy} of the score across thresholds. Bottom: scores
  while \ifoccupied{vacant (gray) and occupied (orange), each scaled to its
  own peak}{vacant, scaled to their peak}. The dotted line marks
  $\tau = \text{\threshold}$.}
  \label{fig:threshold}
\end{figure}

\begin{table}[tb]
  \caption{Readings by observed and predicted state at
  $\tau = \text{\threshold}$.}
  \label{tab:confusion}
  \scriptsize%
  \centering%
  \begin{tabular}{@{}l r r@{}}
    \toprule
    & \multicolumn{2}{c@{}}{Predicted} \\
    \cmidrule(l){2-3}
    Observed & Occupied & Vacant \\
    \midrule
    Occupied & \marimonum{model_summary.true_positive}
      & \marimonum{model_summary.false_negative} \\
    Vacant & \marimonum{model_summary.false_positive}
      & \marimonum{model_summary.true_negative} \\
    \bottomrule
  \end{tabular}%
\end{table}

At $\tau = \text{\threshold}$, the score labels
\marimonum[share]{model_summary.accuracy} of the readings correctly
(\cref{tab:confusion}).
\ifoccupied{%
  Of the readings it calls occupied, \marimonum[share]{model_summary.precision}
  are, and it finds \marimonum[share]{model_summary.recall} of the occupied
  readings. \Cref{fig:threshold} shows the tradeoff behind that choice: a
  lower threshold finds every occupied reading at the cost of false alarms,
  and a higher one loses recall quickly once it passes the scores that
  occupied readings reach. Over the same thresholds, light alone labels at
  most \marimonum[share]{score_baselines.light} of the readings correctly and
  \cotwo{} alone \marimonum[share]{score_baselines.co2}, while the score
  reaches \marimonum[share]{score_baselines.score}.%
}{%
  Every reading it calls occupied is a false alarm, and there are
  \marimonum{model_summary.false_positive} of them. Without occupied readings,
  recall is undefined, and so is precision wherever the score calls a reading
  occupied, so \cref{fig:threshold} plots accuracy alone.%
}

\begin{table}[tb]
  \caption{The five longest episodes of consecutive misjudged readings at
  $\tau = \text{\threshold}$, with their mean light and \cotwo{}.}
  \label{tab:errors}
  \scriptsize%
  \centering%
  \begin{tabular}{@{}l l r r r@{}}
    \toprule
    & & & Light & \cotwo{} \\
    Episode & Outcome & Readings & (lux) & (ppm) \\
    \midrule
    \marimorows[5]{error_episodes}{%
      \marimodate[shortday]{#1.start}, \marimotime{#1.start}%
      \ifnum\marimovalue[1]{#1.readings}>1 --\marimotime{#1.end}\fi
      & \marimovalue{#1.outcome}
      & \marimonum{#1.readings}
      & \marimonum[count]{#1.light}
      & \marimonum[count]{#1.co2} \\}
    \bottomrule
  \end{tabular}%
\end{table}

\Cref{tab:errors} groups the misjudged readings into episodes of consecutive
minutes, longest first. Each row carries the light and \cotwo{} that produced
the error, so its cause can be read against \cref{eq:score}: a false alarm with
bright light and low \cotwo{} was called by light, and a missed reading with
dim light and rising \cotwo{} was not yet reached by \cotwo{}.

\section{Discussion}
\label{sec:discussion}

The score is evaluated on the readings that set its normalization, so its
numbers describe how well it separates these readings in sample, not how it
would do on the next week. The data cover one room for about a week, which is
enough to show the method and too little to choose a threshold for a
building. \ifoccupied{In these readings light alone already separates the
states almost as well as the score, so the weight on \cotwo{} buys little
accuracy here. }{}The two sensors fail in different places, though: daylight
raises light without anyone present, and \cotwo{} lags behind arrivals and
departures. A longer record would show whether keeping both terms pays off.

\section{Conclusion}

\ifoccupied{%
  A score with two named weights reads Room~01's occupancy from light and
  \cotwo{} with the accuracy this paper reports, and every term in it can be
  checked against a chart.%
}{%
  Without occupied readings, the score can only be judged by its false
  alarms, and every one of them can be checked against a chart.%
}

\section*{Supplemental Materials}
\label{sec:supplemental_materials}

The notebook \texttt{examples/occupancy.py}, this paper's \LaTeX{} source in
\texttt{examples/\_\_marimo\_\_/studio/occupancy/paper}, and the notebook's
other views are in the marimo Studio repository at
\url{https://github.com/marimo-team/marimo-studio}.

\section*{Figure Credits}
\label{sec:figure_credits}

The notebook draws every figure with Matplotlib~\cite{Hunter2007Matplotlib}
from the UCI Occupancy Detection data, which is licensed under
\href{https://creativecommons.org/licenses/by/4.0/}{CC BY 4.0}.

\acknowledgments{%
  The occupancy readings were recorded by Luis M. Candanedo and V\'{e}ronique
  Feldheim.%
}

\bibliographystyle{abbrv-doi-hyperref}
\bibliography{references}

\end{document}
